% from ega0-20-fr.tex
\oldpage[0\textsubscript{IV-1}]{131}
Le premier membre de \sref{0.20.5.5.1-fr} n'est autre alors que
$\Omega^1_{B/A}\otimes_A A'$. On peut écrire
$B'\otimes_{A'}B'=(B\otimes_A B)\otimes_A A'$ à un isomorphisme canonique
près et $p_{B'/A'}$ s'identifie alors à
$p_{B/A}\otimes\mathrm{I}_{A'}$ ; par suite (comme $p_{B/A}$ est surjectif)
$\mathfrak{J}_{B'/A'}=\operatorname{Im}(\mathfrak{J}_{B/A}\otimes_A A')$
et
$\mathfrak{J}^2_{B'/A'}
=\operatorname{Im}(\mathfrak{J}^2_{B/A}\otimes_A A')$, d'où
$\mathrm{P}^1_{B'/A'}=\mathrm{P}^1_{B/A}\otimes_A A'$ à un isomorphisme
canonique près, qui transforme en eux-mêmes les idéaux d'augmentation ;
comme le $A$-module $\mathrm{P}^1_{B/A}$ s'identifie canoniquement à la
somme directe de $B$ et de $\Omega^1_{B/A}$, on a
\begin{equation}
\tag{20.5.5.2}
\label{0.20.5.5.2-fr}
  \Omega^1_{B/A}\otimes_A A'=\Omega^1_{B'/A'}
\end{equation}
par le même isomorphisme, et on vérifie aussitôt que le composé de cet
isomorphisme et de l'isomorphisme canonique
$\Omega^1_{B/A}\otimes_B B'\xrightarrow{\sim}
\Omega^1_{B/A}\otimes_A A'$ n'est autre que \sref{0.20.5.5.1-fr}.
\end{proof}

\begin{env}[20.5.6]
\label{0.20.5.6-fr}
Les homomorphismes canoniques \sref{0.20.5.2.4-fr} et
\sref{0.20.5.3.3-fr} donnent, par fonctorialité, pour tout $C$-module $L$,
des homomorphismes canoniques
\begin{align}
\tag{20.5.6.1}
\label{0.20.5.6.1-fr}
  &\operatorname{Hom}_C(\Omega^1_{C/A},L)
  \to\operatorname{Hom}_C(\Omega^1_{B/A}\otimes_B C,L)
  =\operatorname{Hom}_B(\Omega^1_{B/A},L) \\
\tag{20.5.6.2}
\label{0.20.5.6.2-fr}
  &\operatorname{Hom}_C(\Omega^1_{C/B},L)
  \to\operatorname{Hom}_C(\Omega^1_{C/A},L).
\end{align}
Compte tenu de \sref{0.20.4.8.2-fr} et des diagrammes commutatifs
\sref{0.20.5.2.6-fr} et \sref{0.20.5.3.5-fr}, ces homomorphismes ne sont
autres (à une identification canonique près) que les homomorphismes
\sref{0.20.2.1.2-fr} et \sref{0.20.2.1.3-fr} respectivement.
\end{env}

\begin{theorem}[20.5.7]
\label{0.20.5.7-fr}
Soient $u:A\to B$, $v:B\to C$ deux homomorphismes d'anneaux.
\begin{enumerate}
  \item[\rm(i)] La suite de $C$-modules
\begin{equation}
\tag{20.5.7.1}
\label{0.20.5.7.1-fr}
  \Omega^1_{B/A}\otimes_B C
  \xrightarrow{v_{C/B/A}}\Omega^1_{C/A}
  \xrightarrow{u_{C/B/A}}\Omega^1_{C/B}\to0
\end{equation}
est exacte.
  \item[\rm(ii)] Pour que $v_{C/B/A}$ soit inversible à gauche, il faut et
  il suffit que $C$ soit une $B$-algèbre formellement lisse relativement à
  $A$ (pour les topologies discrètes (cf. \sref{0.19.9.1-fr})) ; en
  particulier, il suffit pour cela que $C$ soit une $B$-algèbre formellement
  lisse (pour les topologies discrètes).
\end{enumerate}
\end{theorem}

\begin{proof}
(i) L'exactitude de la suite \sref{0.20.2.4.2-fr} montre tout d'abord,
compte tenu de \sref{0.20.5.6-fr}, que la suite
\[
  0\to\operatorname{Hom}_C(\Omega^1_{C/B},L)
  \to\operatorname{Hom}_C(\Omega^1_{C/A},L)
  \to\operatorname{Hom}_C(\Omega^1_{B/A}\otimes_B C,L)
\]
est exacte pour tout $C$-module $L$. On sait que cela implique l'exactitude
de la suite \sref{0.20.5.7.1-fr} (Bourbaki, \emph{Alg.}, chap.~II,
3\textsuperscript{e} éd., §~2, n\textsuperscript{o}~1, th.~1).

(ii) En vertu de l'exactitude de \sref{0.20.5.7.1-fr}, dire que
$v_{C/B/A}$ est inversible à gauche signifie que la suite
\begin{equation}
\tag{20.5.7.2}
\label{0.20.5.7.2-fr}
  0\to\Omega^1_{B/A}\otimes_B C
  \xrightarrow{v_{C/B/A}}\Omega^1_{C/A}
  \xrightarrow{u_{C/B/A}}\Omega^1_{C/B}\to0
\end{equation}
est \emph{exacte et scindée} ; on sait (Bourbaki, \emph{loc. cit.},
n\textsuperscript{o}~1, prop.~1) que cela équivaut à dire que pour tout
$C$-module $L$, la suite
\[
  0\to\operatorname{Hom}_C(\Omega^1_{C/B},L)
  \to\operatorname{Hom}_C(\Omega^1_{C/A},L)
  \to\operatorname{Hom}_C(\Omega^1_{B/A}\otimes_B C,L)\to0
\]
est exacte ; compte tenu de \sref{0.20.5.6-fr} et
\sref{0.20.2.4.2-fr}, cette condition équivaut à
\[
  \operatorname{Exalcom}_{B/A}(C,L)=0
\]
pour tout $C$-module $L$, et la conclusion résulte donc de
\sref{0.19.9.8.1-fr}.
\end{proof}

